\begin{lemma}
\label{lemma-closed-push-pseudo-coherent}
Let $i : Z \to X$ be a morphism of ringed spaces such that
$i$ is a closed immersion of underlying topological spaces and such that
$i_*\mathcal{O}_Z$ is pseudo-coherent as an $\mathcal{O}_X$-module.
Let $E \in D(\mathcal{O}_Z)$. Then $E$ is $m$-pseudo-coherent
if and only if $Ri_*E$ is $m$-pseudo-coherent.
\end{lemma}

\begin{proof}
Throughout this proof we will use that $i_*$ is an exact functor, and
hence that $Ri_* = i_*$, see Modules, Lemma \ref{modules-lemma-i-star-exact}.

\medskip\noindent
Assume $E$ is $m$-pseudo-coherent. Let $x \in X$. We will find a neighbourhood
of $x$ such that $i_*E$ is $m$-pseudo-coherent on it. If $x \not \in Z$
then this is clear. Thus we may assume $x \in Z$. We will use
that $U \cap Z$ for $x \in U \subset X$ open form a fundamental system of
neighbourhoods of $x$ in $Z$. After shrinking $X$ we may assume $E$ is
bounded above. We will argue by induction on
the largest integer $p$ such that $H^p(E)$ is nonzero. If $p < m$, then
there is nothing to prove. If $p \geq m$, then $H^p(E)$ is an
$\mathcal{O}_Z$-module of finite type, see
Cohomology, Lemma \ref{cohomology-lemma-finite-cohomology}.
Thus we may choose, after shrinking $X$, a map
$\mathcal{O}_Z^{\oplus n}[-p] \to E$ which induces a surjection
$\mathcal{O}_Z^{\oplus n} \to H^p(E)$. Choose a distinguished triangle
$$
\mathcal{O}_Z^{\oplus n}[-p] \to E \to C \to \mathcal{O}_Z^{\oplus n}[-p + 1]
$$
We see that $H^j(C) = 0$ for $j \geq p$ and that $C$ is $m$-pseudo-coherent
by Cohomology, Lemma \ref{cohomology-lemma-cone-pseudo-coherent}.
By induction we see that $i_*C$ is $m$-pseudo-coherent on $X$.
Since $i_*\mathcal{O}_Z$ is $m$-pseudo-coherent on $X$ as well, we conclude
from the distinguished triangle
$$
i_*\mathcal{O}_Z^{\oplus n}[-p] \to i_*E \to i_*C \to
i_*\mathcal{O}_Z^{\oplus n}[-p + 1]
$$
and 
Cohomology, Lemma \ref{cohomology-lemma-cone-pseudo-coherent}
that $i_*E$ is $m$-pseudo-coherent.

\medskip\noindent
Assume that $i_*E$ is $m$-pseudo-coherent. Let $z \in Z$.
We will find a neighbourhood of $z$ such that $E$
is $m$-pseudo-coherent on it. We will use
that $U \cap Z$ for $z \in U \subset X$ open form a fundamental system of
neighbourhoods of $z$ in $Z$. After shrinking $X$ we may assume $i_*E$
and hence $E$ is bounded above. We will argue by induction on
the largest integer $p$ such that $H^p(E)$ is nonzero. If $p < m$, then
there is nothing to prove. If $p \geq m$, then $H^p(i_*E) = i_*H^p(E)$
is an $\mathcal{O}_X$-module of finite type, see
Cohomology, Lemma \ref{cohomology-lemma-finite-cohomology}.
Choose a complex $\mathcal{E}^\bullet$ of $\mathcal{O}_Z$-modules
representing $E$. We may choose, after shrinking $X$,
a map $\alpha : \mathcal{O}_X^{\oplus n}[-p] \to i_*\mathcal{E}^\bullet$
which induces a surjection
$\mathcal{O}_X^{\oplus n} \to i_*H^p(\mathcal{E}^\bullet)$.
By adjunction we find a map
$\alpha : \mathcal{O}_Z^{\oplus n}[-p] \to \mathcal{E}^\bullet$
which induces a surjection
$\mathcal{O}_Z^{\oplus n} \to H^p(\mathcal{E}^\bullet)$.
Choose a distinguished triangle
$$
\mathcal{O}_Z^{\oplus n}[-p] \to E \to C \to \mathcal{O}_Z^{\oplus n}[-p + 1]
$$
We see that $H^j(C) = 0$ for $j \geq p$. From the distinguished triangle
$$
i_*\mathcal{O}_Z^{\oplus n}[-p] \to i_*E \to i_*C \to
i_*\mathcal{O}_Z^{\oplus n}[-p + 1]
$$
the fact that $i_*\mathcal{O}_Z$ is pseudo-coherent
and 
Cohomology, Lemma \ref{cohomology-lemma-cone-pseudo-coherent}
we conclude that $i_*C$ is $m$-pseudo-coherent.
By induction we conclude that $C$ is $m$-pseudo-coherent.
By Cohomology, Lemma \ref{cohomology-lemma-cone-pseudo-coherent}
again we conclude that $E$ is $m$-pseudo-coherent.
\end{proof}